\subsection{Redemption-Cap Design: The Ex-Ante Arithmetic}\label{sec:discussion-capdesign}

That the caps would sit above achievable runoff is not my observation. New York Fed staff projected \$20--\$30 billion in monthly principal payments when the \$35 billion ceiling took effect \citep{nyfed2022}. Their May 2022 baseline path implied agency MBS runoff of \$740.6 billion over the window, against the caps' \$1{,}417.5 billion. What this paper adds is the quantification, constructed ex ante. The \$35 billion cap was set roughly 1.7 to 1.9 times the Federal Reserve's own contemporaneous projection of achievable runoff; the multiple depends on the same disclosed intra-2022 allocation choice as the anticipated share (Section~\ref{sec:method-benchmark}). From the other side, both components of the $\beta_1 = 0$ null---scheduled amortization from the book's coupon, age, and term composition, and baseline turnover at any assumed floor---recover 88.7\% of the realized shortfall at the in-sample calibration point and 85.7\% at the off-window floor under the production floor form (35.6\% under the additive form). Table~\ref{tab:wal} prices the book's rate-inelastic principal path in maturity units. Under the empirical CPR the book's approximate weighted-average life is 9.4 years against 3.4 years at 2021 realized speeds. That 3.4-year comparator is a refinancing-boom baseline, not a turnover norm. So the speed that would have filled the cap was the one the window's rate configuration had already shut off.

\begin{table}[H]
\centering
\small
\begin{threeparttable}
\caption{Approximate SOMA portfolio WAL by prepayment scenario (years).}
\label{tab:wal}
\begin{tabular}{@{}lrr@{}}
\toprule
Scenario (mean CPR) & WAL, June 2022 & WAL, Nov.\ 2025 \\
\midrule
Scheduled amortization only (0\%) & 14.7 & 12.6 \\
Empirical path, ABM-basis back-out (5.14\%) & 9.4 & 8.5 \\
Empirical path, note-rate WAC basis (5.79\%) & 8.9 & 8.1 \\
Production ABM (11.68\%) & 6.0 & 5.6 \\
Hazard Path A (3.34\%) & 10.9 & 9.7 \\
Hazard Path B (4.76\%) & 9.7 & 8.7 \\
Danish rule-only leg, production anchor (5.61\%) & 9.1 & 8.2 \\
Danish-level bracketing leg (3.39\%) & 10.8 & 9.6 \\
Baseline turnover floor, off-window headline read (4.99\%) & 9.5 & 8.6 \\
No-shock counterfactual, 2021 realized speeds (22.81\%) & 3.4 & 3.3 \\
\bottomrule
\end{tabular}
\begin{tablenotes}[flushleft]\footnotesize
\item Single-pool-per-cell approximation: vintage face shares from the SOMA CUSIP tabulation (2021: 43.9\%; 2020: 16.3\%; 2017--19: 6.0\%; 2022: 23.1\%; pre-2017: 10.6\%), a 90.7\%/9.1\% 30-year/15-year term split, coupon at the book's 2.49\% WAC, and each scenario's window-mean CPR held constant. Level-payment amortization with monthly SMM; November 2025 values age the same cells 42 months. Cell WAL formula: for remaining term $n$ months, $\mathrm{WAL} = \sum_{t=1}^{n} t\, p_t / (12 \sum_{t} p_t)$, with $p_t$ month-$t$ principal (scheduled level-payment principal plus $\mathrm{SMM} = 1-(1-\mathrm{CPR})^{1/12}$ of the surviving balance). Worked anchors for the scheduled-only row: a 2.49\% 30-year pool has origination WAL 16.9 years (the 2021-vintage cell, aged 12 months at June 2022, is at 16.3); the 9.1\% 15-year sleeve (origination WAL 8.0 years) and the older vintages (32.9\% of face is 2020 or earlier) pull the book blend to 14.7. Representative cell ages ($\{2022{:}\,2;\ 2021{:}\,12;\ 2020{:}\,32;\ 2017\text{--}19{:}\,60;\ \text{pre-}2017{:}\,120\}$ months at June 2022) are single-pool stand-ins recovered by matching the printed table, and anchored externally: ages derived from the SOMA CUSIP tabulation's own back-derived origination dates (maturity minus term, face-weighted within cells, computed with no reference to this table) run $\{2021{:}\,12;\ 2020{:}\,28;\ 2017\text{--}19{:}\,50\}$ months and give a scheduled-only blend of 14.9 years; naive vintage midpoints likewise give 14.9. The printed 14.7 sits within 0.3 years of both alternatives. Coverage reconciliation: the 90.7\%/9.1\% split (which sums to the 99.8\% covered) is the audit parse of the \$1{,}940.8 billion book (the 2026-07-15 as-of of the same object whose committed June-2022 as-of total is \$2{,}700.6 billion; the difference is the interval's runoff, not a discrepancy), with 0.2\% of face in other terms never modeled; the 90.6\% of Appendix~\ref{sec:robustness-pipeline} is the production run's own as-of-date 30-year-only coverage share; the two parses differ by as-of date and their components round independently.
\item Accounting basis of each row. The two empirical rows are the same SOMA series backed out under two coupon conventions: 5.14\% is the ABM-basis back-out, 5.79\% the note-rate WAC --- the corrected amortization basis, since the note rate is what amortizes the loan --- and the hazard-basis back-out at 5.52\% sits between them at 9.1/8.3 years. Moving from the ABM basis to the note-rate basis shortens the empirical path by 0.5 years at the window's open and 0.4 at its close (run \texttt{wal\_note\_rate\_basis}, through the identical calculator), which is the same order as the 0.6-year Danish rule-only effect this table also reports. Every other row is a \emph{simulated} mean CPR from its own leg and carries no back-out convention. The 6.0-year extension is deliberately \emph{not} restated on the note-rate basis: the no-shock row's 22.81\% is a 2021 back-out on the ABM basis with no committed note-rate counterpart, so differencing across the two would mix conventions inside one statistic.
\item Turnover benchmark: the final row is a refinancing-boom baseline rather than a normal-turnover one --- 22.81\% is the 2021 realized speed, most of which is refinancing --- and the turnover-floor row is its normal-turnover counterpart. The paper's own out-of-window turnover measurement is the floor read (Table~\ref{tab:oosfloor}): defensible 2018 reads of 4.70\%, 4.99\% and 5.33\% by depth of out-of-the-moneyness, against refi-contaminated pooled 2017--2019 and 2019 anchors at 6.07\% and 6.91\%. The turnover-floor row prices the middle read --- the paper's headline floor --- through the identical calculator (run \texttt{wal\_normal\_turnover}): it falls between the Path B and empirical ABM-basis rows, and monotonicity in CPR puts the band ends on either side of it.
\end{tablenotes}
\end{threeparttable}
\end{table}

The design question that follows is speculative: design space, not a finding these estimates support. Whether this episode generalizes, this design can pose precisely but not answer. \citet{du2024}'s cross-country record covers seven central banks running active and passive runoff, and the mechanism here is not specific to the Federal Reserve: a fixed nominal cap above the rate-inelastic principal path of a book whose coupons sit below market. What would settle it is narrow and public: each bank's announced cap or redemption schedule against its realized redemptions, month by month, over its own tightening window. I do not assert the other six were non-binding; I record that the comparison is constructible from published data and is not made here. Pre-committing to a substitution instrument---a standing policy of substituting Treasury runoff or outright MBS sales when passive runoff undershoots---would convert extension risk from a silent shortfall into an explicit, priced policy choice, and leaves the cap free to communicate the path, not force it. Indexing the cap to the observable state of the outstanding stock is the weaker alternative: it still works through the instrument the shortfall defeats.

What the decomposition supplies is computable in advance of the cycle, as a band: scheduled amortization is fixed by the book's composition; the turnover floor is measured, its defensible out-of-window reads carrying the central-elasticity marginal from $+6.8$ points at the deepest out-of-the-moneyness cut to $+4.3$ at the shallowest (Table~\ref{tab:oosfloor}). The floor read's own sampling error widens the identified margin further: the mid-cut read's binding wild-cluster inversion runs from 4.033\% to 6.181\% (the Webb rung's narrower read, 4.177\% to 5.800\%, is the one the cap grid below was run at) (Section~\ref{sec:robustness-floor}). Propagated through the committed floor-to-marginal grid, that sampling error prices the lock-in marginal at \$0.42 to \$1.66 billion per month of foregone passive principal. That band attaches to the identified marginal, not to the two projection levels, which are accounting constructions with no sampling interval of their own. On my committed projections, over the 42 window months the phased ceilings averaged \$33.75 billion per month of allowable MBS redemption against the passive principal the book could deliver, \$17.6 billion per month under the uniform-spread allocation and \$20.0 billion per month under the settlement-aware one---ratios of 1.91 and 1.69.

The indexing proposal can now be priced: against the window-minimum market rate, the share of book face sitting below market is 99.9\% at a zero out-of-the-moneyness cut and 99.6\% at cuts of 25 and 50 basis points, the last two coinciding because the book's coupons sit on a half-point grid. Carrying each cut's committed floor read through the identical amortization gives achievable passive principal, hence an implied cap, of \$17.07, \$16.47 and \$15.95 billion per month (run \texttt{state\_contingent\_cap\_grid}). The grid supplies the missing function: the whole depth range moves the implied cap by about \$1.1 billion per month, while the mid-cut floor read's own wild-cluster interval moves it by about \$2.9 billion, from \$15.02 to \$17.88 billion per month---the observable is near-degenerate, and dominated by the measurement error in the floor it would be indexed to. One scope limit travels with it: the observable is the SOMA book's bucketed coupon distribution while the floor is estimated on the Freddie panel's loan-level gaps---two populations at two grains. A redemption cap paces and communicates the portfolio's composition path, and Section~\ref{sec:method-benchmark} concedes these caps never bound over this window, so a binding level trades the communication objective for a speed the schedule was not built to deliver. A cap intended to bind would sit at or below the lower edge of the book's projected scheduled-plus-turnover principal path, with the identified elasticity margin ($+2.1$ to $+13.2$ points of the shortfall across the calibration box) as the behavioral band around it. That is the form the Federal Reserve's own contemporaneous projection took: a range. Nothing in my estimates identifies the welfare or market effects of such designs; active sales in particular would realize the mark-to-market losses passive runoff defers.

Two boundaries belong beside this design lesson. First, non-bindingness may have been a choice, not an error. A Treasury-symmetric cap is simple to communicate and forecloses forced sales, and nothing in the public record establishes that the Committee intended the MBS cap to bind. Still, the FOMC's stated longer-run intention to hold primarily Treasury securities gives the MBS-specific shortfall a standing the cap alone does not. Against that composition objective, slower runoff is delayed progress whatever the cap's intent. Second, the design space is priced in partial equilibrium. A cap set to bind, or active sales, would transmit to primary-market spreads, and through them to the lock-in marginal itself. Pricing that feedback is outside this paper's accounting.

\paragraph{What a designer of a future runoff program should and should not take.} Three quantities transfer: the baseline turnover floor in CPR units (Section~\ref{sec:robustness-floor}), the scheduled-amortization path (computable from the book at purchase time), and the elasticity band as the behavioral budget. Two numbers should not travel alone: the $+5.6$-point anchor, a convention inside a range rather than a central tendency, and the \$764.7 billion cap-shortfall, a statement about a never-binding schedule's composition rather than a missed objective. One accounting rule: face value answers the benchmark question, cash the reserve-drain question, and the two differ by the par windfall---a transfer, not a resource cost (Appendix~\ref{app:params}).
